\documentclass[12pt,a4paper]{article}
\usepackage[T1]{fontenc}
\usepackage[utf8]{inputenc}
\usepackage{amsmath,amssymb}
\usepackage{siunitx}
\usepackage{geometry}
\geometry{margin=2.5cm}

\title{Cloud Physics \& Precipitation -- Sample Solutions}
\author{}
\date{}

\begin{document}
\maketitle

\section*{Q1 -- Latent Heating in an Updraft}

Given:
\[
w = \SI{20}{\meter\per\second}, \quad
\Delta q_v = 0.005 \,\si{\kilo\gram\per\kilo\gram}, \quad
L = 2.5\times 10^6 \,\si{\joule\per\kilo\gram}, \quad
c_p \approx 1004 \,\si{\joule\per\kilo\gram\per\kelvin}.
\]

Released latent heat per unit mass:
\[
\Delta Q = L \Delta q_v
= 2.5\times 10^6 \times 0.005
= 1.25\times 10^4 \,\si{\joule\per\kilo\gram}.
\]

Temperature increase:
\[
\Delta T = \frac{\Delta Q}{c_p}
\simeq \frac{1.25\times 10^4}{1004}
\approx \SI{12.4}{\kelvin}.
\]

If the condensation occurs over \SI{1}{\kilo\meter}, the time is
\[
\Delta t = \frac{\Delta z}{w}
= \frac{1000}{20}
= \SI{50}{\second},
\]
so the heating rate is
\[
\frac{\Delta T}{\Delta t} \approx
\frac{12.4}{50} \approx \SI{0.25}{\kelvin\per\second}.
\]

\section*{Q2 -- Evaporative Cooling and Negative Buoyancy}

Given:
\[
\Delta q_v = +0.008 \,\si{\kilo\gram\per\kilo\gram}.
\]

Evaporative cooling:
\[
\Delta T = -\frac{L \Delta q_v}{c_p}
= -\frac{2.5\times 10^6 \times 0.008}{1004}
\approx -\SI{20}{\kelvin}.
\]

Buoyancy (taking a typical environmental temperature
\(T \approx \SI{300}{\kelvin}\)):
\[
B = g\,\frac{T' }{T_{\text{env}}}
\approx 9.81 \,\frac{-20}{300}
\approx -0.65 \,\si{\meter\per\second\squared}.
\]

\section*{Q3 -- Why \texorpdfstring{$\Gamma_m < \Gamma_d$}{Γm < Γd}?}

Dry adiabatic lapse rate:
\[
\Gamma_d = \frac{g}{c_p} \approx
\SI{9.8}{\kelvin\per\kilo\meter},
\]
arises from adiabatic expansion and cooling of unsaturated air.

In saturated air, condensation releases latent heat, partially
offsetting adiabatic cooling. Thus the temperature decreases more
slowly with height:
\[
\Gamma_m \approx \SI{6}{\kelvin\per\kilo\meter}
< \Gamma_d.
\]

\section*{Q4 -- CAPE with/without Virtual Temperature Correction}

Parcel is warmer than environment by about
\(\Delta T \approx \SI{8}{\kelvin}\) from the LCL up to \SI{500}{\hecto\pascal}
(rough approximation).

Buoyancy without virtual correction:
\[
B \approx g \frac{\Delta T}{T}
\approx 9.81 \,\frac{8}{265}
\approx 0.30 \,\si{\meter\per\second\squared},
\]
assuming a representative \(T \approx \SI{265}{\kelvin}\).

If the LFC is around \(z \approx \SI{0.8}{\kilo\meter}\) and
\SI{500}{\hecto\pascal} at \(z \approx \SI{5}{\kilo\meter}\),
then \(\Delta z \approx \SI{4.2}{\kilo\meter}\).

Approximate CAPE:
\[
\text{CAPE} \approx \overline{B}\,\Delta z
\approx 0.30 \times 4200
\approx 1.3\times 10^3 \,\si{\joule\per\kilo\gram}.
\]

Including virtual temperature correction with
\(\Delta T_v \approx \SI{7}{\kelvin}\):
\[
B_v \approx 9.81 \,\frac{7}{265}
\approx 0.26 \,\si{\meter\per\second\squared},
\]
\[
\text{CAPE}_v \approx 0.26 \times 4200
\approx 1.1\times 10^3 \,\si{\joule\per\kilo\gram}.
\]

\section*{Q5 -- CIN (Convective Inhibition)}

Definition:
\[
\text{CIN} = -\int_{z_s}^{z_{\text{LFC}}}
g\,\frac{T_{v,p} - T_{v,e}}{T_{v,e}}\,dz.
\]

If the environment is on average warmer than the parcel
below the LFC (negative buoyancy), one may approximate
\[
\text{CIN} \approx |\overline{B}_{\text{neg}}|\,\Delta z.
\]

For a required lift of \(\Delta p \approx \SI{50}{\hecto\pascal}\),
the corresponding depth is roughly
\(\Delta z \sim 400\text{--}\SI{500}{\meter}\).
Assuming \(|\overline{B}| \approx 0.1\,\si{\meter\per\second\squared}\):
\[
\text{CIN} \sim 0.1 \times 450
\approx \SI{45}{\joule\per\kilo\gram}.
\]

\section*{Q6 -- Supersaturation Estimate}

Given:
\[
z \approx \SI{10}{\kilo\meter}, \quad
w = \SI{25}{\meter\per\second}, \quad
\tau_{\text{cond}} = \SI{200}{\second}.
\]

Take a characteristic vertical distance
\(\Delta z \approx \SI{1}{\kilo\meter}\):
\[
\tau_{\text{asc}} \approx \frac{\Delta z}{w}
= \frac{1000}{25} = \SI{40}{\second}.
\]

A simple box model suggests:
\[
S \sim \frac{\tau_{\text{asc}}}{\tau_{\text{cond}}}
= \frac{40}{200} = 0.2
\;\Rightarrow\;
S \approx 20\% .
\]

\section*{Q7 -- Boussinesq Approximation and Equations}

Assume:
\begin{itemize}
  \item \(\rho = \rho_0(z) + \rho'\), with \(|\rho'| \ll \rho_0\),
  \item \(\rho_0\) depends only on height \(z\),
  \item flow is nearly incompressible: \(\nabla\cdot\mathbf{v} \approx 0\),
  \item density variations only appear in the buoyancy term.
\end{itemize}

Start with the vertical momentum equation:
\[
\frac{Dw}{Dt} = -\frac{1}{\rho}\frac{\partial p}{\partial z} - g + \nu \nabla^2 w.
\]

Decompose:
\[
p = p_0(z) + p', \quad
\rho = \rho_0(z) + \rho',
\]
with basic state hydrostatic balance:
\[
\frac{\partial p_0}{\partial z} = -\rho_0 g.
\]

Subtract the basic state:
\[
\frac{Dw}{Dt} =
-\frac{1}{\rho_0}\frac{\partial p'}{\partial z}
+ g\frac{\rho'}{\rho_0}
+ \nu \nabla^2 w.
\]

Relate density and potential temperature perturbations:
\[
\frac{\rho'}{\rho_0}
\approx -\frac{\theta'}{\theta_0},
\quad
B = g\frac{\theta'}{\theta_0}.
\]

Hence the Boussinesq momentum and continuity equations:
\begin{align*}
\frac{D\mathbf{v}}{Dt}
&= -\frac{1}{\rho_0}\nabla p'
   + B\,\hat{\mathbf{k}}
   + \nu \nabla^2\mathbf{v},\\
\nabla\cdot\mathbf{v} &= 0.
\end{align*}

\section*{Q8 -- Moist Adiabatic Lapse Rate}

One standard result for the moist adiabatic lapse rate is:
\[
\Gamma_m
= \Gamma_d
\frac{1 + \dfrac{L q_v}{R_d T}}
     {1 + \dfrac{L^2 q_v}{c_p R_v T^2}},
\]
where
\(\Gamma_d = g/c_p\),
\(q_v\) is the (saturated) mixing ratio, \(R_d\) is the gas constant
for dry air, and \(R_v\) for water vapor.

Sketch of derivation:
\begin{itemize}
  \item Use the first law of thermodynamics for a saturated parcel
        with condensation:
        \[
        c_p dT - L \,dq_s = -\alpha\,dp,
        \]
        with \(q_s(T,p)\) the saturation mixing ratio.
  \item Combine with hydrostatic balance
        \(\mathrm{d}p/\mathrm{d}z = -\rho g\) and the ideal gas law.
  \item Express \(dq_s\) in terms of \(dT\) and \(dp\).
  \item Manipulate algebraically to obtain the expression above.
\end{itemize}

\section*{Q9 -- Max Updraft from CAPE}

Given CAPE:
\[
\text{CAPE} = \SI{3000}{\joule\per\kilo\gram}.
\]

Parcel theory approximation:
\[
w_{\max} \approx \sqrt{2\,\text{CAPE}}
= \sqrt{2\times 3000}
= \sqrt{6000}
\approx \SI{77.5}{\meter\per\second}.
\]

\section*{Q10 -- Tilting of Streamwise Vorticity}

Given:
\[
\eta = 0.015 \,\si{\per\second},
\quad
S = 0.02 \,\si{\per\second}.
\]

Tilting contribution to vertical vorticity:
\[
\left(\frac{d\omega}{dt}\right)_{\text{tilt}}
\sim \eta S
= 0.015 \times 0.02
= 3\times 10^{-4} \,\si{\per\second\squared}.
\]

Over a characteristic time
\(t \sim \SI{100}{\second}\):
\[
\omega \sim
\left(\frac{d\omega}{dt}\right)t
\approx 3\times 10^{-4} \times 100
= 0.03 \,\si{\per\second}.
\]

\section*{Q11 -- Negative Buoyancy and Surface Winds}

Given:
\[
T_p = -5^\circ\text{C} = \SI{268}{\kelvin},\quad
T_e = 2^\circ\text{C} = \SI{275}{\kelvin}.
\]

Temperature difference:
\[
\Delta T = T_p - T_e = -7\,\si{\kelvin}.
\]

Buoyancy:
\[
B = g\,\frac{T_p - T_e}{T_e}
\approx 9.81 \frac{-7}{275}
\approx -0.25 \,\si{\meter\per\second\squared}.
\]

If the descending cold air has depth \(\Delta z\):
\[
\text{DCAPE} \sim |B|\,\Delta z.
\]
For \(\Delta z \approx \SI{3}{\kilo\meter}\):
\[
\text{DCAPE} \sim 0.25\times 3000
\approx 750 \,\si{\joule\per\kilo\gram}.
\]
Maximum downdraft speed:
\[
w_{\max} \sim \sqrt{2\,\text{DCAPE}}
\approx \sqrt{1500}
\approx \SI{39}{\meter\per\second}.
\]

Such a downdraft spreading horizontally at the surface can produce
surface winds on the order of \(30\text{--}\SI{40}{\meter\per\second}\).

\section*{Q12 -- Pressure Perturbation from Vorticity}

Given:
\[
\omega = 0.04 \,\si{\per\second},\quad
R = \SI{2}{\kilo\meter} = 2000\,\si{\meter}.
\]

For solid-body rotation,
\[
\zeta = 2\Omega = \omega
\quad\Rightarrow\quad
\Omega = \frac{\omega}{2} = 0.02 \,\si{\per\second}.
\]

Tangential velocity at radius \(R\):
\[
v(R) = \Omega R = 0.02 \times 2000
= \SI{40}{\meter\per\second}.
\]

Cyclostrophic balance:
\[
\frac{v^2}{R}
= \frac{1}{\rho}\frac{\partial p'}{\partial r}.
\]

For solid-body rotation:
\[
\frac{\partial p'}{\partial r}
= \rho\,\Omega^2 r.
\]

Pressure difference from center to radius \(R\):
\[
\Delta p = p(0) - p(R)
= \int_0^R \rho \Omega^2 r \, dr
= \frac{1}{2}\rho\Omega^2 R^2.
\]

Assume \(\rho \approx \SI{1}{\kilo\gram\per\meter^3}\):
\[
\Delta p = \frac{1}{2}\times 1 \times
(0.02)^2 \times (2000)^2
= 0.5 \times 4\times 10^{-4} \times 4\times 10^6
= 800\,\si{\pascal}
\approx 8\,\si{\hecto\pascal}.
\]

\section*{Q13 -- Why the Right-Mover Deviates from the Mean Wind}

In an environment with strong vertical wind shear and positive
storm-relative helicity, the updraft tilts horizontal vorticity
into the vertical, creating a mesocyclone (cyclonic vertical vorticity).

The associated dynamic pressure perturbation (low pressure in the
updraft core) pulls the storm toward the right of the shear vector,
so the right-moving supercell deviates significantly from the
mean wind.

\section*{Q14 -- Straight vs Curved Hodograph and Supercells}

\begin{itemize}
  \item \textbf{Straight-line hodograph:} shear is symmetric; splitting storms
        can produce both right- and left-movers, often with similar potential
        intensity.
  \item \textbf{Clockwise-curved hodograph (NH):} storm-relative helicity is
        strongly positive for right-movers, favoring long-lived, dominant
        cyclonic supercells.
\end{itemize}

Thus, a clockwise-curved hodograph is most favorable for
classical, long-lived right-moving supercells.

\section*{Q15 -- Curvature and Saturation Vapor Pressure (Kelvin Effect)}

Surface tension causes a pressure difference across a curved
liquid interface:
\[
p_\text{inside} - p_\text{outside} = \frac{2\sigma}{r},
\]
where \(\sigma\) is surface tension and \(r\) the droplet radius.

This raises the chemical potential of the liquid, so
the equilibrium vapor pressure over a curved surface
is greater than over a flat surface. The Kelvin equation:
\[
e_s(r) = e_s(\infty)\,
\exp\left(\frac{2\sigma}{\rho_w R_v T}\frac{1}{r}\right),
\]
or
\[
\ln\frac{e_s(r)}{e_s(\infty)}
= \frac{2\sigma}{\rho_w R_v T}\,\frac{1}{r}.
\]

As \(r\) decreases, \(e_s(r)\) increases.

\section*{Q16 -- Köhler Theory for a Tiny NaCl Particle (Concept)}

Köhler equation (one common form):
\[
S(r) = \frac{e(r)}{e_s(\infty)}
= a_w(r)\,\exp\left(\frac{A}{r}\right),
\]
which is often written as
\[
\ln S = \frac{A}{r} - \frac{B}{r^3},
\]
where
\[
A = \frac{2\sigma}{\rho_w R_v T}
\]
(Kelvin term), and
\[
B \propto n_s
\]
(Raoult term, with \(n_s\) the moles of solute; for NaCl,
van 't Hoff factor \(i \approx 2\)).

The critical radius and supersaturation:
\[
\frac{d(\ln S)}{dr} = 0
\quad\Rightarrow\quad
-\frac{A}{r_c^2}
+ \frac{3B}{r_c^4} = 0
\quad\Rightarrow\quad
r_c = \left(\frac{3B}{A}\right)^{1/2},
\]
\[
S_c = \exp\left(\frac{A}{r_c} - \frac{B}{r_c^3}\right).
\]

For an extremely small NaCl mass (here
\(\sim 10^{-16}\,\si{\gram}\)), \(B\) is very small;
thus \(r_c\) is near the dry radius and
\(S_c\) is very large (supersaturation well above a few percent,
possibly $>100\%$). Such a tiny particle is not an efficient CCN
in typical atmospheric conditions.

\section*{Q17 -- Bergeron Process at \texorpdfstring{$-15^\circ$C}{-15°C}}

At \(-15^\circ\text{C}\), the saturation vapor pressure over water
\(e_{sw}\) is larger than that over ice \(e_{si}\). Typical values:
\[
e_{sw}(-15^\circ\text{C}) \approx 1.9 \,\si{\hecto\pascal},\quad
e_{si}(-15^\circ\text{C}) \approx 1.4 \,\si{\hecto\pascal},
\]
so
\[
\frac{e_{sw}}{e_{si}} \sim 1.3.
\]

If the environment is saturated with respect to liquid water
(\(e \approx e_{sw}\)), then it is supersaturated with respect to ice:
\[
S_i \approx \frac{e}{e_{si}}
\approx \frac{e_{sw}}{e_{si}}
\sim 1.3,
\]
i.e.\ about 30\,\% supersaturated relative to ice.

Thus, ice crystals grow rapidly by vapor deposition, while liquid
droplets tend to evaporate: the Bergeron process favors ice growth.

\section*{Q18 -- CAPE vs Dynamic Pressure in a Supercell Updraft}

\begin{itemize}
  \item \textbf{Buoyancy (CAPE):} provides thermal energy; large CAPE
        gives a large potential max updraft
        \(w_{\max} \sim \sqrt{2\,\text{CAPE}}\).
  \item \textbf{Dynamic pressure perturbations:} arise from strong
        vertical shear and mesocyclone rotation. Negative pressure
        perturbations in the updraft core enhance vertical acceleration
        and can be comparable to or even exceed buoyancy locally.
\end{itemize}

In many strong supercells, the updraft is maintained by
a combination of buoyancy and dynamic (pressure-gradient)
forcing; in the updraft core, dynamic pressure forces can
be as important as buoyancy, or sometimes dominant.

\end{document}
